% ID: SEL-D0016-BA-E02
% ESTADO: cotejada. Segunda pasada visual completa y borrador compilado revisados.
% Fuente: Selectividad/matematicas_aplicadas/matematicas_aplicadas_P_D.pdf
% SHA-256: 11865b83e509d9cc2723c752f7a81f714378e1b651f4a4d010dab4a4576b74c2
% Página: 1; bloque A; ejercicio original 2.

Se consideran las matrices
$A=\begin{pmatrix}1&0&-2\\a&1&0\end{pmatrix}$,
$B=A\cdot A^t$ y $C=\begin{pmatrix}1&-2\\1&0\end{pmatrix}$,
siendo $a$ un parámetro real.
\begin{enumerate}[label=\alph*)]
  \item \textbf{(0.75 puntos)} \textquestiondown{}Para qué valores del
  parámetro $a$ existe la inversa de la matriz $B$?
  \item \textbf{(0.75 puntos)} Para $a=1$, calcule la inversa de la matriz $B$.
  \item \textbf{(1 punto)} Para $a=1$, resuelva la ecuación matricial
  $B^t\cdot X+9\,C=O$.
\end{enumerate}
